\documentclass[output=paper]{langscibook}
\author{Authorname\orcid{}\affiliation{}}
\title{Title}
\abstract{Abstract}
\IfFileExists{../localcommands.tex}{
  \addbibresource{../localbibliography.bib}
  \usepackage{tabularx,multicol}
\usepackage[hyphens]{url}
\urlstyle{same}

\usepackage{langsci-optional}
\usepackage{langsci-lgr}
\usepackage{langsci-gb4e}

\usepackage{soul}
% \usepackage{booktabs}
% \usepackage{geometry}
\usepackage{multirow}
% \usepackage{makecell} %safely breaks a line inside a table cell
% \usepackage{tablefootnote} %footnotes in table captions
\usepackage{enumerate}
\usepackage{tikz}
\usepackage{array}
\usepackage[shortlabels]{enumitem}

%for having several figures on one line, turning words in 90°angles
% \usepackage{graphicx}
\usepackage{subcaption}

% \usepackage{caption} %for making the caption line longer in concrete cases

% \usepackage{rotating} %rotate a whole page; for large tables

\usepackage{fontspec}
\newfontfamily\charis{Charis SIL} %for being able to use the symbol ‿ and a better command of some of the boldfaced diacritics in 01
\newfontfamily\japanesefont{Noto Serif CJK JP} % for Japanese script
% \usepackage{tipa}

% \usepackage{needspace} %to keep exs on one page

\usepackage{xcolor}

\usepackage{pgfplots}
\pgfplotsset{compat=1.18}

%%%%%   AVMs for 02	%%%%%
\usepackage{langsci-avm}
\avmsetup{switch = ?}


\usepackage{langsci-branding}

  %\newcommand*{\orcid}{}


\makeatletter
\let\thetitle\@title
\let\theauthor\@author
\makeatother

\newcommand{\togglepaper}[1][0]{
   \bibliography{../localbibliography}
   \papernote{\scriptsize\normalfont
     \theauthor.
     \titleTemp.
     To appear in:
     E. Di Tor \& Herr Rausgeberin (ed.).
     Booktitle in localcommands.tex.
     Berlin: Language Science Press. [preliminary page numbering]
   }
   \pagenumbering{roman}
   \setcounter{chapter}{#1}
   \addtocounter{chapter}{-1}
}

\newbool{bookcompile}
\booltrue{bookcompile}
\newcommand{\bookorchapter}[2]{\ifbool{bookcompile}{#1}{#2}}



\renewcommand{\thempfootnote}{\fnsymbol{mpfootnote}} % use * for footnotes in float env

% \newcolumntype{L}[1]{>{\raggedright\arraybackslash}p{#1}} %left-aligned (non-justified) text in tables
% \newcolumntype{R}[1]{>{\raggedleft\arraybackslash}p{#1}} % right-aligned (non-justified) text in tables

\definecolor{customgreen}{RGB}{27,158,119}
\definecolor{custompink}{RGB}{231,42,138}
\definecolor{customblue}{RGB}{117,112,179}
\definecolor{customlightblue}{RGB}{143,170,220}
\definecolor{customdarkblue}{RGB}{68,114,196}
\definecolor{customyellow}{RGB}{225,192,0}
\definecolor{customorange}{RGB}{233,137,21}
\definecolor{customgray}{RGB}{229,229,229}


% Left-aligned indexes
% % \makeatletter
% % \patchcmd{\theindex}{\twocolumn}{\raggedright\twocolumn}{}{}
% % \makeatother

%to make Japanese scripts in the bibliography smaller
\newcommand{\japbib}[1]{{\japanesefont\small #1}}

%%%%%%%%%%%%%%%%%%%	MACROS for 02	%%%%%%%%%%%%%%%%

%%% Formatting	%%%%%

\newcommand{\textex}[1]{\textit{#1}} 		 %for formatting linguistic examples in-text%
%command-shift X

\newcommand{\lxm}[1]{\textsc{#1}}  		%for formatting names of lexemes
%command-option-shift L

%%%%%%%%% Abbreviations	%%%%%%

\newcommand{\featname}[1]{\textsc{#1}}
\newcommand{\featspec}[2]{\featname{#1}:{#2}}
\newcommand{\feat}[2]{\textsc{#1}:{#2}}  				% \feat{FEATURE}{value}  sc's
\newcommand{\featnameonly}[1]{\textsc{#1}} 			% \featnameonly{featurename}  sc's

\newcommand{\pr}{$^{\prime}$}

\newcommand{\Corr}{\textit{Corr}}
\newcommand{\propm}{\textit{pm}}

\newcommand{\PC}{Π$_{C}$}

\newcommand{\PF}{Π$_{F}$}
\newcommand{\PR}{Π$_{R}$}

\newcommand{\lex}{$\mathcal{L}$}
 
  %% hyphenation points for line breaks
%% Normally, automatic hyphenation in LaTeX is very good
%% If a word is mis-hyphenated, add it to this file
%%
%% add information to TeX file before \begin{document} with:
%% %% hyphenation points for line breaks
%% Normally, automatic hyphenation in LaTeX is very good
%% If a word is mis-hyphenated, add it to this file
%%
%% add information to TeX file before \begin{document} with:
%% %% hyphenation points for line breaks
%% Normally, automatic hyphenation in LaTeX is very good
%% If a word is mis-hyphenated, add it to this file
%%
%% add information to TeX file before \begin{document} with:
%% \include{localhyphenation}
\hyphenation{
    Al-ex-an-der
    Ber-ber
    chan-ges
    cir-cum-stan-ces
    coin-age
    Fraj-zyngier
    ita-liano
    Ja-pa-nese
    Kö-nig
    Krzy-żanowski
    li-te-ra-tur-no-go
    mi-zen-kei
    par-a-digm
    phe-nom-e-non
    ren-tai-shi
    Slo-vene
    Spen-cer
    Stoj-kov
    Ver-bi-ca-re-se
    wide-spread
}

\hyphenation{
    Al-ex-an-der
    Ber-ber
    chan-ges
    cir-cum-stan-ces
    coin-age
    Fraj-zyngier
    ita-liano
    Ja-pa-nese
    Kö-nig
    Krzy-żanowski
    li-te-ra-tur-no-go
    mi-zen-kei
    par-a-digm
    phe-nom-e-non
    ren-tai-shi
    Slo-vene
    Spen-cer
    Stoj-kov
    Ver-bi-ca-re-se
    wide-spread
}

\hyphenation{
    Al-ex-an-der
    Ber-ber
    chan-ges
    cir-cum-stan-ces
    coin-age
    Fraj-zyngier
    ita-liano
    Ja-pa-nese
    Kö-nig
    Krzy-żanowski
    li-te-ra-tur-no-go
    mi-zen-kei
    par-a-digm
    phe-nom-e-non
    ren-tai-shi
    Slo-vene
    Spen-cer
    Stoj-kov
    Ver-bi-ca-re-se
    wide-spread
}
 
  \togglepaper[1]%%chapternumber
}{}

\begin{document}
\maketitle 
%\shorttitlerunninghead{}%%use this for an abridged title in the page headers
% ATTENTION: Diacritics on the following phonetic characters might have been lost during conversion: {'ə', 'ɛ'}

\title{The conditions of uninflectability in nouns in the Slavic languages}%L\footnote{ \textrm{I am very much obliged to my reviewers for their comments which have significantly contributed to improving this chapter. All flaws are of course mine.}}}

Ursula Doleschal

\textit{In Slavic languages, major word classes like nouns, adjectives, and verbs typically inflect for various grammatical categories. However, some lexemes within these classes are uninflected, meaning they do not display any grammatical markers. I first address how to define uninflectedness, which is straightforward for nouns but more complex for adjectives and verbs. Next, I explore which lexeme classes are most commonly uninflected in languages such as Czech, Polish, Russian, Serbian, Slovene, and Slovak, focusing primarily on proper names, loanwords, and abbreviations, which are mostly nouns. I highlight that uninflectedness exists on a continuum, with different degrees of uninflectedness within each class, and this continuum is inversely related to other factors like numeral and compound formation. I argue that the lack of integration of certain nouns into the inflectional system is linked to specific declension features like stem-ending, base form, and grammatical gender. These features vary in restrictiveness across languages, with some (e.g., Russian) showing higher uninflectedness in nouns, while others (e.g., Slovene) are more permissive. Additionally, noun classes with feminine gender tend to be more restrictive in terms of stem-structure and base-form than those with masculine gender.}

\section{1Introduction}

In this article I will present some findings of my habilitation thesis \citep{Doleschal2000} where I investigated the phenomenon of uninflectedness and uninflectability in six Slavic languages – Russian, Polish, Slovak, Czech, Slovene and Serbian. The aim of this investigation was to find out how and why in highly inflecting languages, both single words such as the Russian noun \textit{metro} as in \textit{stancija} \textit{metro} ‘metro station’, the Czech adjective \textit{blond} in \textit{blond vlasy} ‘blond hair’or the numeral \textit{dve} in Serbian as in \textit{s dve devojke} ‘with two girls’ and entire classes of words as Russian and Polish nouns ending in the vowel /u/, e.g., \textit{s kenguru} ‘with (a) kangaroo’ remain partly or even fully uninflected, and eventually, interjections as \textit{pryg} ‚hop‘ in Russian which might qualify for uninflectable verbs: \textit{my pryg} ‘we hopped’.

In the course of this research, other phenomena also related to uninflectedness emerged, most notably constructions with company names in compound-like structures, such as \textit{Duracell} in \textit{Duracell nagradna igra} (Slovene) ‘Duracell prize draw’, \textit{Milka} in \textit{50 Milka plišanih kravica} (Serbian) ‘50 Milka fluffy toy cows’ or a historical growth of inflectedness as wittnessed by the (relatively) new inflectional class \textit{kuli} in Czech, Slovak and Polish. These issues have been addressed in Doleschal (1999; 2007; 2014; 2015).

The structure of the present article is the following: Firstly, the Slavic languages will be characterised as inflecting languages (\sectref{sec:key:2}).

Secondly, the different forms of the phenomenon which have been identified in the material will be described (\sectref{sec:key:3}). 

Finally, the ways in which uninflectability manifests itself in the six languages under investigation, and how this uninflectability is related to the inflectional classes in Slavic, will be examined, as will the research question about the conditions of uninflectedness (\sectref{sec:key:4}).

\section{Slavic languages as inflecting languages}  %2 /

The Slavic languages are traditionally classified as highly inflecting and belonging to the fusional type (cf. e.g., \citealt{ComrieCorbett1993}: 6). This means that they signal grammatical meanings by endings which are attached to lexical morphemes, usually stems. Endings are characterised by cumulative exponence, i.e., one ending expresses more than one grammatical meaning (feature value). Inflection is different and characteristic for the different parts of speech.

The inflected parts of speech are: noun, pronoun, adjective, numeral, determiner, verb (including participle; Sussex, \citealt{Cubberley2006}: 222), and – depending on one’s theoretical premises – adverb (cf. \citealt{Nagórko2009}: 743). Within the inflected parts of speech, word forms are organised in paradigms forming inflectional classes so that any lexeme of these parts of speech necessarily belongs to (usually) one such class (cf. \tabref{tab:key:1}). The grammatical (morphosyntactic and morphosemantic) categories involved in inflection are the following:\footnote{ \textrm{Cf. \citet{SpencerWiemer2020} for a discussion of inflectional categories in Slavic languages.}} case, number, gender, animacy, (gradation) in nouns and adjectives; person, number, gender, tense, mood, voice, (aspect); gerund, participle in verbs, determining the content paradigm of a lexeme which belongs to a certain part of speech. For example, the Polish nouns \textit{pan} ‘gentleman, mister’ and \textit{żona} ‘woman’ take different sets of inflectional endings for the same values of the categories case and number (\tabref{tab:key:1}). Furthermore, these two nouns are classified differently with regard to the category of gender – while \textit{pan} has the feature masculine, \textit{żona} displays feminine gender. Whether the gender feature is also expressed by the sets of endings along with the case and number features is controversial (cf. \citealt{Doleschal2009}: 148). 


\begin{tabularx}{\textwidth}{XXXXX}
 & {\bfseries Sg.} & {\bfseries Pl.} & {\bfseries Sg.} & {\bfseries Pl.}\\
\lsptoprule
N & pan-Ø & pan-owie & żon-a & żon-y\\
G & pan-a & pan-ów & żon-y & żon\\
D & pan-u & pan-om & żon-ie & żon-om\\
A & pan-a & pan-ów & żon-ę & żon-y\\
I & pan-em & pan-ami & żon-ą & żon-ami\\
L & pan-u & pan-ach & żon-ie & żon-ach\\
V & pani-e & pan-owie & żon-o & żon-y\\
\lspbottomrule
\end{tabularx}
%%please move \begin{table} just above \\begin{tabular . 
\begin{table}
\caption{Inflectional classes (Pol.): pan (M) ‚gentleman, mister‘, żona (F)‚woman‘}
\label{tab:key:1}
\end{table}

\section{Types and manifestations of uninflectedness} %3 /
\subsection{General} %3.1 /

In the general area of uninflectedness, different manifestations of the phenomenon can be discerned:

1) classes of words that never take any endings (3.2)

2) single words that never take any endings (3.2)

3) words that may or may not take endings and vary by syntactic context, register, sociolect, speaker (or freely) (3.3).

\subsection{Uninflectability (paradigmatic uninflectedness)} %3.2 /

In the first place, there are words that never take any inflectional endings and are therefore paradigmatically uninflected, as the borrowings \textit{miss} f. ‘miss’, e.g., in \textit{miss Polonia} or \textit{gnu} n., f., m. (in Russian, Polish, Slovak) ‘gnu’ for a type of antelope. Paradigmatic uninflectedness pertains to the lexical level. It is an inherent property of the lexeme. This manifestation of uninflectedness will be called “uninflectability”.

Uninflectable lexemes can be subdivided into two groups: 

1. In the first group uninflectability is regular, i.e., it can be described by rules based on narrowly linguistic (e.g., phonological or lexical) properties: in Russian, e.g., nouns ending in the vowel /u/ cannot be inflected: \textit{tabu} ‘taboo’, \textit{ragu} ‘ragout’, \textit{kenguru} ‘kangaroo’, \textit{Lulu} “Loulou”, \textit{Subaru} “Subaru”. This case shall be called “systematic uninflectability”.

2. The second group contains words whose uninflectability cannot be explained by properties of the language system, as in the case of some toponyms in Slovak ending in the vowel /e/ as \textit{Skopje} or in a so-called soft consonant\footnote{\textrm{In the Slavic languages, consonants underwent several palatalisations, leading to correlations between soft (palatal or palatalised) and hard (non-palatal) consonants. This difference is widely morphologised, so that “hard” stems and “soft” stems select different sets of endings which are not conditioned by phonology. Moreover, there are “neutral” sc. “unpaired” consonants which do not take part in the correlation.}} (cf. \citealt{ComrieCorbett1993}: 6; 9-10) as \textit{Providence} (/s/) while \textit{Opole}, \textit{Bystrice} (both ending in /e/), \textit{Provence} (ending phonologically in /s/) with the same type of phonological ending are inflectable. This case shall be called “unsystematic uninflectability”.

On this basis we can now define what exactly will be understood by ‘uninflectability’:

A word is called uninflectable,\footnote{\textrm{Cf. also} \textrm{Krzyżanowski 2013: 24-35 for a thorough terminological discussion.}} if

1) it belongs to an inflectable part of speech, and

2) does not take any endings normally found with words of that kind, but

3) may nevertheless appear in any syntactic context typical for the part of speech in question.

Condition 1 is necessary to distinguish the words in question from uninflected parts of speech such as conjunctions and prepositions in the Slavic languages. 

Condition 2 distinguishes uninflectable words from semi-inflectable ones as, e.g., Polish \textit{muzeum} ‘museum’ (\tabref{tab:key:2}), which is uninflected in the singular but regularly inflected in the plural, and from contextual uninflectedness (see below 3.3).


\begin{tabularx}{\textwidth}{XXX}

\lsptoprule

\textbf{~} & \textbf{Sg.} & \textbf{Pl.}\\
\textbf{N} & muzeum & muze-a\\
\textbf{G} & muzeum & muze-ów\\
\textbf{D} & muzeum & muze-om\\
\textbf{A} & muzeum & muze-a\\
\textbf{I} & muzeum & muze-ami\\
\textbf{L} & muzeum & muze-ach\\
\textbf{V} & muzeum & muze-a\\
\lspbottomrule
\end{tabularx}
%%please move \begin{table} just above \\begin{tabular . 
\begin{table}
\caption{semi-inflectable nouns in Polish}
\label{tab:key:2}
\end{table}

Condition 3 is necessary to distinguish between uninflectability proper and defectivity (cf. \citealt{Doleschal2000}, \citealt{Spencer2020}): for example, the Russian indefinite numeral \textit{malo} ‘a little, few’ has only one form and appears solely in Nominative and Accusative Singular expressions \REF{ex:key:1}: 

\ea%1
    \label{ex:key:1}
    \gll \\
         \\
    \glt
    \z % you might need an extra \z if this is the last of several subexamples

         a. Tam malo detej 

          there few.NOM children.GEN’

         ‘There are few children there’

b. Ja znaju malo detej 

    ‘I know few.ACC children.GEN’ 

    ‘I know few children

c. *Ona interesuetsja malo detej 

      she is-interested-in few.INSTR children 

      ‘She is interested in few children. 

It is thus also uninflectable, but not in the same way as are nouns as \textit{kenguru} ‘kangaroo’ which as such can express all syntactic cases and both singular and plural, e.g. \REF{ex:key:2}: 

\ea%2
    \label{ex:key:2}
    \gll \\
         \\
    \glt
    \z % you might need an extra \z if this is the last of several subexamples

         a. Tam molod-oj kenguru  

         There young-M.NOM.SG kangaroo.NOM.SG

         ‘There is a young kangaroo there’

b.    Ja tam vižu kenguru 

      ‘I there see kangaroo.ACC.SG/PL’

      ‘I see a kangaroo there’ / ‘I see kangaroos there’, 

c.    Ona interesuetsja kenguru 

      she is-interested-in cangaroo.INSTR.SG/PL

      ‘She is interested in kangaroos / a kangaroo’. 

This definition can be exemplified with two female names taken from Czech – \textit{Jana} and \textit{Mimi} (cf. \tabref{tab:key:3}). As proper names, both are clearly nouns and thus fulfil condition 1. Secondly, it can be seen from \tabref{tab:key:3} that \textit{Mimi} does not take any case endings but occurs with the same case meanings as the inflected lexeme \textit{Jana}, i.e., it has the same content paradigm as \textit{Jana}, and thus fulfils condition 2. It can therefore appear in all necessary syntactic contexts as is illustrated by the use of prepositions in the second column and so also fulfils condition 3.


\begin{tabularx}{\textwidth}{XXXX}

\lsptoprule

Nom. &  & Jan-a & Mimi\\
Gen. & (bez) & Jan-y & Mimi\\
Dat. & (k) & Jan{}-ě & Mimi\\
Acc. & (pro) & Jan-u & Mimi\\
Loc. & (při) & Jan-ě & Mimi\\
Instr. & (s) & Jan-ou & Mimi\\
Voc. &  & Jan-o! & Mimi!\\
\lspbottomrule
\end{tabularx}
%%please move \begin{table} just above \\begin{tabular . 
\begin{table}
\caption{inflectable and uninflectable name in Czech}
\label{tab:key:3}
\end{table}

Condition 3 also implies that an uninflectable word possesses and represents all morphological categories and their grammatical meanings (feature values) of the part of speech in question. But how large can such a paradigm be?  Would it also have to cover otherwise periphrastic morphology? Does it have to cover word-class changing morphology? 

The Russian verb paradigm has up to 124 forms (e.g., the verb \textit{ljubit’} ‘to love’), including the four participles with their paradigms and the two gerunds (regarding aspect as a lexical feature, cf. \citealt{Genis2021}).

Does an uninflectable item have to substitute for all of these forms in order to be called uninflectable and not defective? For my material, I have come to the following conclusion both on an empirical and a theoretical basis: an uninflectable lexeme must be able to independently represent all synthetic forms of the paradigm of a part of speech in question, for the nominal parts of speech this means all the case-number forms in the noun and case-number-gender forms and adverb in the adjective, and case forms in the numeral as well as case-gender forms for some numerals, but not the periphrastic forms, such as periphrastic gradation.

This is mainly due to the fact that words that are usually perceived as uninflectable can combine with periphrastic markers such as Czech \textit{více} in \textit{více blond} ‘more blond’ for the comparative or as Russian \textit{pryg by} ‘would hop’ for the conditional. From the theoretical point of view, it has been argued (cf. e.g., \citealt{Haspelmath2000}) that obligatory preposition-noun combinations, such as in the locative in the Slavic languages, are usually seen as phrases although they could just as well be analysed as periphrastic case forms. However, we do not expect an uninflectable noun in the (syntactic) locative to co-signal a preposition. Another viewpoint could be that periphrasis is on the borderline between syntax and inflection \citep{BrownEtAl2012}, hence not a prototypical inflectional marker and could therefore be neglected. As to the word-class changing forms (adverb, participle), \citet{Spencer2020} has argued for excluding them from the paradigm, and I will follow him here.

Under these premises an uninflectable adjective (as \textit{super}) would substitute for the following paradigm (\tabref{tab:key:4}) in Slovene.


\begin{tabularx}{\textwidth}{XXXXXXX}
 & m. & f. & n. & m. & f. & n.\\
\lsptoprule
Sg. & ~ & ~ & ~ & ~ & ~ & ~\\
Nom. & nov-Ø/nov-i\footnotemark{} & nov-a & nov-o & super & super & super\\
Gen. & nov-ega & nov-e & nov-ega & super & super & super\\
Dat. & nov-emu & nov-i & nov-emu & super & super & super\\
Acc. & nov-Ø/nov-ega\footnotemark{} & nov-o & nov-o & super & super & super\\
Loc. & nov-em & nov-i & nov-em & super & super & super\\
Instr. & nov-im & nov-o & nov-im & super & super & super\\
Du. & ~ & ~ & ~ & ~ & ~ & ~\\
Nom. & nov-a & nov-i & nov-i & super & super & super\\
Gen. & nov-ih & nov-ih & nov-ih & super & super & super\\
Dat. & nov-im & nov-im & nov-im & super & super & super\\
Acc. & nov-a & nov-i & nov-i & super & super & super\\
Loc. & nov-ih & nov-ih & nov-ih & super & super & super\\
Instr. & nov-imi & nov-imi & nov-imi & super & super & super\\
Pl. &  &  &  &  &  & \\
Nom. & nov-i & nov-e & nov-a & super & super & super\\
Gen. & nov-ih & nov-ih & nov-ih & super & super & super\\
Dat. & nov-im & nov-im & nov-im & super & super & super\\
Acc. & nov-e & nov-e & nov-a & super & super & super\\
Loc. & nov-ih & nov-ih & nov-ih & super & super & super\\
Instr. & nov-imi & nov-imi & nov-imi & super & super & super\\
\lspbottomrule
\end{tabularx}
\addtocounter{footnote}{-2}
\stepcounter{footnote}\footnotetext{ \textrm{indefinite/definite form (cf. \citealt{Greenberg2008}: 36)}}
\stepcounter{footnote}\footnotetext{ \textrm{inanimate/animate gender}}
%%please move \begin{table} just above \\begin{tabular . 
\begin{table}
\caption{inflectable and uninflectable adjective in Slovene}
\label{tab:key:4}
\end{table}

Since gradation can be periphrastic in all Slavic languages, this category need not be expressed by an indeclinable adjective by itself. In Slovene, adjectives can form synthetic forms of comparative and superlative as \textit{nov-ejš-i} 'new-GRADATION-M.SG' 'newer' and \textit{naj-nov-ejš-i} SUP-new-GRADATION-M.SG' 'newest' along with a periphrastic comparative and superlative, which is then also applied to uninflectable adjectives (\tabref{tab:key:5}).


\begin{tabularx}{\textwidth}{XXXX}
 & \textbf{Positive} & \textbf{Comparative} & \textbf{Superlative}\\
\lsptoprule
synthetic & nov-Ø, -a, -o & nov-ejš-i, -a, -e & naj-nov-ejš-i, -a, -e\\
periphrastic & nov-Ø, -a, -o

super & bolj nov-i, -a, -o

bolj super & najbolj nov-i, -a, -o

najbolj super\\
\lspbottomrule
\end{tabularx}
%%please move \begin{table} just above \\begin{tabular . 
\begin{table}
\caption{Gradation of adjectives in Slovene} 
\label{tab:key:5}
\end{table}

\subsection{3.2 Contextual (syntagmatic)\footnotemark{} uninflectedness}
\footnotetext{ \textrm{In \citet{Doleschal2000} “syntagmatic” was used throughout. However, as one of the reviewers pointed out, this term is too narrow. Therefore, the phenomenon is called “contextual uninflectedness” here.}}

Having defined uninflectability, let us now turn to a closely related phenomenon which is often treated together with uninflectability under the same heading (“nesklonjaemost'”, “neohebnost”) but is distinct from the first: contextual uninflectedness. Contextual uninflectedness means that a word, which is fully inflectable on the lexical level, remains uninflected for contextual reasons, such as text type, style, sociolect and last but not least – syntactic context, e.g., apposition or use of a classifier as in \textit{v naselj-u} \textbf{\textit{Njivice}} (in village-LOC.SG Njivice‘in the village Njivice’) in Slovene or \textit{w województw-ie Białystok} (in province-LOC.SG ‘in the province Białystok’. However, such uninflectedness is usually not obligatory, e.g., in Polish, masculine surnames ending in /e/ or /i/ are inflectable, but especially in collocation with a first name or a title often remain uninflected, e.g., \textit{Maksimilian-a Kolbe} Maksimilian-GEN.SG Kolbe [GEN.SG] instead of \textit{Maksimilian-a Kolb-ego} Maksimilian-GEN.SG Kolbe-GEN.SG (cf., e.g., Krzyżanowski 2013: 39-40). In Slovene this happens with feminine surnames in /a/, e.g., \textit{z Ad-o Muha} ‘with Ada-INSTR.SG Muha [INSTR.SG]’ ‘with Ada Muha’ instead of \textit{z Ad-o Muh-o ‘}Muha{}-INSTR.SG’. In Serbian, male surnames can remain uninflected in collocation with first names, if they precede the first name: \textit{Petrović Ivan-a} Petrović.NOM.SG Ivan-GEN.SG ‘Petrović Ivan’s’.

On the other hand, in Czech abbreviations as \textit{NATO} or \textit{FIFA} remain uninflected in official language, but are inflected in casual speech (Internetová jazyková příručka (2008–2024)). Another case in point are Russian place names formed with the suffixes -\textit{ovo} and -\textit{ino} such as \textit{Perovo}, \textit{Puškino} which are in principle inflectable and should be inflected according to the norms of Contemporary Standard Russian (cf. Paxomov n. y., Redakcija portala «Gramota.ru» n.y., \citealt{Zaliznjak2008}: 739, \citealt{Ageenko2010}: 35). Nevertheless, they often remain uninflected both in spoken and written language and seem to show a tendency towards uninflectability (\citealt{Graudina1980}:156-59, \citealt{Panov1968}: 56-59, \citealt{Zaliznjak2008}: 739, but also \citealt{Tejkin2021} for an opposite view). 

One of the interesting issues is if and how contextual and paradigmatic uninflectedness relate to each other, i.e., is contextual uninflectedness a preliminary stage of uninflectability? This may be the case as with the aforementioned Russian place names, but it is certainly not the case, when an uninflected noun occurs together with an appositive classifier in Slovene or in Serbian. This problem needs an investigation of the relationship between the effects of register and style and language change caused by both sociolinguistic and inherently morphologocial factors, such as productivity of inflectional classes, and has to be left for further exploration.

\subsection{Geographical distribution of uninflectability in six Slavic languages} %3.3 /

As a result of the study by \citet{Doleschal2000} the manifestation of uninflectability in the six Slavic languages under investigation can be characterised as follows: Uninflectability occurs regularly, but in different ways, with the parts of speech \textit{noun}, \textit{adjective}, \textit{(cardinal)} \textit{numeral}. The degree of uninflectability manifests itself as a kind of continuum with two trajectories for different parts of speech: one from west to east and the other from south to north (cf. \tabref{tab:key:6}). This means that for nouns the highest degree of uninflectability is found in Russian and the lowest in Slovene, whereas for both adjectives and numerals the highest degree is found in Serbian and the lowest (if any) in Russian.


\begin{tabularx}{\textwidth}{XXXXXX}

\lsptoprule
\multicolumn{1}{c}{} & \multicolumn{1}{c}{\textbf{Noun}} & \multicolumn{1}{c}{} & \multicolumn{1}{c}{} & \multicolumn{1}{c}{} & \\
%\hhline%%replace by cmidrule{~----~}
N & Polish &  &  & Russian & N\\
%\hhline%%replace by cmidrule{~----~}
${\uparrow}$ & Czech &  &  & Slovak & ${\uparrow}$\\
%\hhline%%replace by cmidrule{~----~}
S & Slovene &  &  & Serbian & S\\
%\hhline%%replace by cmidrule{~----~}
\multicolumn{1}{c}{} & \multicolumn{1}{c}{W} & \multicolumn{1}{c}{${\rightarrow}$} & \multicolumn{1}{c}{\raggedleft ${\rightarrow}$} & \multicolumn{1}{c}{\raggedleft E} & \\
\lspbottomrule
\end{tabularx}

\begin{tabularx}{\textwidth}{XXXXXX}

\lsptoprule
\multicolumn{1}{c}{} & \multicolumn{1}{c}{\textbf{Numeral}} & \multicolumn{1}{c}{(only cardinal numeral)} & \multicolumn{1}{c}{} & \multicolumn{1}{c}{} & \\
%\hhline%%replace by cmidrule{~----~}
N & Polish &  &  & Russian & N\\
%\hhline%%replace by cmidrule{~----~}
${\downarrow}$ & Czech &  &  & Slovak & ${\downarrow}$\\
%\hhline%%replace by cmidrule{~----~}
S & Slovene &  &  & Serbian & S\\
%\hhline%%replace by cmidrule{~----~}
\multicolumn{1}{c}{} & \multicolumn{1}{c}{W} & \multicolumn{1}{c}{${\rightarrow}$} & \multicolumn{1}{c}{\raggedleft ${\rightarrow}$} & \multicolumn{1}{c}{\raggedleft E} & \\
\lspbottomrule
\end{tabularx}

\begin{tabularx}{\textwidth}{XXXXXX}

\lsptoprule
\multicolumn{1}{c}{} & \multicolumn{1}{c}{\textbf{Adjective}} & \multicolumn{1}{c}{} & \multicolumn{1}{c}{} & \multicolumn{1}{c}{} & \\
%\hhline%%replace by cmidrule{~----~}
N & Polish &  &  & Russian & N\\
%\hhline%%replace by cmidrule{~----~}
${\downarrow}$ & Czech &  &  & Slovak & ${\downarrow}$\\
%\hhline%%replace by cmidrule{~----~}
S & Slovene &  &  & Serbian & S\\
%\hhline%%replace by cmidrule{~----~}
\multicolumn{1}{c}{} & \multicolumn{1}{c}{W} & \multicolumn{1}{c}{${\rightarrow}$} & \multicolumn{1}{c}{\raggedleft ${\rightarrow}$} & \multicolumn{1}{c}{\raggedleft E} & \\
\lspbottomrule
\end{tabularx}
\begin{stylelsTableHeading}
\textmd{\textit{\tabref{tab:key:6}: Characterisation of uninflectability in six Slavic languages}}
\end{stylelsTableHeading}

In \tabref{tab:key:6}, dark slots indicate a high degree of uninflectability, while light slots mean a low degree. A white slot indicates total lack of uninflectability for a certain part of speech, as is, e.g., the case for cardinal numerals in Russian, Polish and Czech where all members of this class display overt inflectional forms. In Serbian, on the other hand, numerals from ‘five’ to ‘999’ are all uninflectable. A black slot thus means that a considerable number of lexical items is affected by uninflectability. 

The shading does not correspond to absolute numbers, but rather to a qualitative criterion: how many lexical classes are affected? How many types of words are affected by non-assignment and how regular is this? In Russian, all classes of names, foreign words and abbreviations are affected by uninflectability. In addition, nouns with any phonological ending and any gender are affected, but systematically those ending in /e/, /i/ and /u/. Most varieties of abbreviations are also uninflectable. Thus, uninflectability of nouns also shows a high degree of regularity or systematicity. This is why the cell for the noun in \tabref{tab:key:6} is shaded black for Russian. In Slovene, on the other hand, names or abbreviations. Re not affected by uninflectability at all. Moreover, there is no phonological ending that would preclude nouns beforehand from being inflected. There are some uninflectable nouns in Slovene, but they cannot be described by rules. Therefore, the corresponding cell is shaded in light grey.

The conditions for this uninflectability are different for the different parts of speech: In numerals we find a diachronic reduction of inflection over the past 100-150 years (\citealt{Suprun1969}, \citealt{Doleschal2000}: 297-313). In this case we are dealing with contextual uninflectedness as a step towards uninflectability. 

Uninflectable adjectives, on the other hand, are almost always borrowings (cf. \citealt{Doleschal2000}: 277-297). Therefore, the reason for uninflectability in adjectives can be found in their lack of integration and their foreignness.

In the case of nouns, the situation is more complicated and will therefore be analysed in detail in the following section.

\section{Uninflectability in nouns} %4 /
\subsection{Introductory considerations}  %4.1 /

Uninflectable nouns can be divided into three lexically and morphologically defined classes, which overlap logically to some extent: \textit{proper nouns}, \textit{abbreviations} and \textit{borrowings}. These groups may intersect: many of the uninflectable proper nouns are at the same time borrowings or foreign names, e.g., the (Italian) male name \textit{Mario}, or the (German) female name \textit{Sigrid} (in Russ., Pol., Cz., Slk.), the toponyms \textit{Oslo}, \textit{Cetinje} (in Russian), and chrematonyms\footnote{ \textrm{also called pragmatonyms, see \citealt{Zgusta1996}.}} as \textit{Subaru} (in Russ., Pol., Cz., Slk.). It is very seldom the case that genuinely native names are uninflectable, as, e.g., the Russian male first name \textit{Sadko}\footnote{ \textrm{a mythical hero of the bylinas, cf. the opera by Rimsky-Korsakov.}} or Russian surnames as \textit{Dolgix}, \textit{Živago} and the Czech type of surname ending in /u:/ (graphic <ů> ) such as \textit{Janů}, \textit{Petrů} in Slovak.\footnote{ \textrm{Both the Russian and the Czech surnames of these types are frozen Genitive forms. In Czech such surnames may be inflected.}} Abbreviations may of course also be borrowings and/or names (e.g., of organisations, such as \textit{IBM}), but they may also be native, such as \textit{SNG} ‘CIS’ (Russ.), \textit{ORMO} (Ochotnicza Rezerwa Milicji Obywatelskiej ‘Volunteer Reserve of the Citizens' Militia’) (Polish) or compounds as \textit{kompolka} ‘commander of a regiment’ (Russian). The class of borrowings includes (besides names) all borrowed common nouns that are not morphologically integrated for various reasons, e.g. \textit{alibi}, \textit{tabu} (Russ., Pol., Cz., Slk.), \textit{miss} ‘miss’, \textit{ledi} ‘lady’ (all investigated languages). 

Apart from these, there are small subgroups of native uninflectable common nouns in some languages, e.g., female designations of persons (\textit{mami} ‘mummy’ in Slovak and Slovene, \textit{professor} ‘woman-professor’ etc. in Polish\footnote{ \textrm{The Polish cases are, however, defective, in that they cannot be used in the plural *}\textrm{\textit{panie professor} }\textrm{‘Ms:N.PL professor’.}}), extragrammatical\footnote{ \textrm{see \citet{Dressler2000} for a discussion of this term.}} formations, such as \textit{cucu} (Slovak) or \textit{ciuciu} (Polish) ‘candy’, \textit{psipsi} ‘wee-wee’) or the designations of letters and musical notes (\textit{be}, \textit{ce}, \textit{de} ‘b, c, d’; \textit{do} “C”\textit{, re} “D”\textit{, mi} “E” in all investigated languages besides Slovene). Other native common nouns are very rarely uninflectable (cf. e.g., Krzyżanowski 2013: 18). 

In the literature on uninflectability various explanations have been brought forward for motivating the uninflectability of nouns. The main reason is that they do not fit the requirements of inflection, for phonological as well as morphological reasons (cf., e.g., Krzyżanowski 2013: 37). Thus, \citet[225]{Isačenko1954} justifies the uninflectability of borrowings on /o/ in Russian, which in the majority of cases bear final stress, with the lack of a morphonological model among Russian nouns (the accent type \textit{okno} /oˈkno/ ‘window’ comprises only a few Russian native words and is not productive), and \citet{Worth1966} attributes uninflectability in Russian to a vocalic stem ending, as would be the case in \textit{bo-a} ‘boa’. Krzyżanowski (2013: 30) mentions a word-final stressed vowel as one cause of uninflectability of abbreviations in Polish, such as \textit{EKG} [ɛkaˈgɛ] ‘electrocardiogram’. 

A notorious reason is also the mismatch between gender and phonological shape of a noun, as with \textit{miss} ‘miss’, \textit{madam} ‘madam’ in Russian \citep[148]{Timberlake2004}. Another reason given is the recognisability of names, e.g., surnames that are homonymous with common nouns, such as Polish \textit{Lato}, \textit{Musik} (Krzyżanowski 2013: 41) which is especially important in official style. This reason is also mentioned as a motivation for not inflecting Russian place names as (neuter) \textit{Puškino}, which could be confused with the paronymous (masculine) place name \textit{Puškin} due to identical endings in the oblique cases (cf., \citealt{Tejkin2021} for a discussion of this argument). Finally, foreignness is presented as a key factor for uninflectability, cf. the discussion by Krzyżanowski (2013: 35-42), also \citet[1075-1076]{Tejkin2021}

It is also put forward that some nouns are assigned more easily or regularly to inflection than others. The latter either remain uninflectable for a span of time or vary between inflection and uninflectedness. This is, e.g., the case with borrowed common nouns and names ending in /o/ as well as with male names ending in /e/ or /i/ in Polish (cf. Krzyżanowski 2013: 20-21).

But what does the fact that certain nouns do not (easily) fit into inflection imply for linguistic explanation? In this regard, two points can be made:

1) Systematic uninflectability demonstrates the limits and thereby the inner structure of the inflectional system of a language. This is to say that in the case of systematic uninflectability there is a mismatch between some property of the noun in question and the rules conditioning inflection in a specific language.

2) Unsystematic uninflectability is most often a phenomenon of prescriptive norms, but it may also reveal tendencies within inflection.

\subsection{Uninflectability and inflectional classes} %4.2 /

We shall commence with an examination of the phenomenon of unsystematic uninflectability. In Russian, for instance, French names ending in stressed /a/ are uninflectable (as \textit{Nikol’a} [nʲɩkʌˈla] “Nicolas”, \textit{Dega} [Dɛˈga] “Degas”. This is not true for names in stressed /a/ from other languages (e.g., \textit{Abdulla} [ʌbdulˈla] “Abdullah”). It is typical for unsystematic uninflectability that such a norm may change, e.g., \textit{Marseille} (ending in the soft consonant /j/) was uninflectable in Slovak according to the norm in 1998 and is now codified as both inflectable and uninflectable (cf. \citealt{Považaj1998,Považaj2013}). Unsystematic uninflectability may therefore be indicative of a transitory state of a (set of) noun(s) towards inflection. 

In other cases, unsystematic uninflectability may be an indicator of a tendency, as with nouns ending in /o/ in Polish, where uninflectability is gaining ground both following and counter to prescriptive norms (cf. \citealt{Doleschal2000}: 56; 68-70; 98-100; 113; 126; 138 for discussion and references). E.g., male surnames in /o/ should be generally inflectable, but in dictionaries many foreign names are marked as uninflectable, such as \textit{Eco}, \textit{Franco}. Other names in /o/, too, often remain uninflected in usage (cf. \citealt{Doleschal2000}: 68-70, Krzyżanowski 2013: 21, 41). Borrowed common nouns in /o/ with neuter gender, on the other hand, are uninflectable at first, but usually come to be inflected over time, e.g., \textit{agio} ‘agio’ already in the 19th century, \textit{auto} ‘car’, \textit{radio} ‘radio’, \textit{kino} ‘cinema’ in the period before the Second World War, and more recently \textit{molo} ‘pier’, \textit{kakao} ‘coacoa’, which are classified as variably inflected and uninflected (cf. \citealt{Doleschal2000}: 126, Żmigrodzki 2004-2021, \citealt{Woliński2020}, consulted 23.11.2023). The integration of neuters in /o/ into a corresponding inflectional class is therefore not automatic, but slower and not always complete (see Dressler, \citealt{Dziubalska-KołaczykFabiszak1997}: 102, \citealt{Orzechowska1974}: 112). This points to a development towards systematic uninflectability for nouns ending in /o/ in Polish.

Let us now discuss systematic uninflectability as an indicator of the limits of the inflectional system. Systematic uninflectability means that a certain type of noun cannot be assigned to any inflectional class of a given language and thus shows what the requirements of inflectional classes are.

The fundamental question is: what characteristics must a word possess, in order to be integrated into the morphological system of a language, i.e. to be assigned to an inflectional class. One important factor in the Slavic languages is the phonological ending. As previously stated, nouns ending in the vowel /u/ cannot be assigned to any inflectional class in Russian or Polish. This implies that the Russian and Polish inflectional classes do not permit a base form in /u/. Consequently, inflectional classes can presuppose a certain base form type, which is a conditio sine qua non for the assignment of a word to them. This type of property of an inflectional class will be referred to as a “constitutive property” (cf. \citealt{Doleschal1997,Doleschal2006}). 

Usually native words correspond to the requirements of the morphological system, because lexicon and regular word formation of a given language provide words with a permissible base form. Therefore, it is often foreign words that are affected by uninflectability, although this is not the only case: relict forms created with synchronously unproductive means of word formation can also be affected, e.g., the (archaic) Russian first names \textit{Mixajlo}, \textit{Danilo} (cf. \citealt{Zaliznjak2008}: 739). Therefore, the inflectional system may also change in a way that leaves native words outside. Therefore, uninflectability enables us to identify, firstly, the constitutive properties associated with inflectional classes and, secondly, the degree of productivity of the latter (cf. \citealt{Wurzel1984}, \citealt{Dressler1997,Dressler2003}). 

In the literature on Slavic languages, the integration of a noun into an inflectional class is usually related to gender and to the phonological ending of the base form as well as to the ending of the stem (\citealt{HentschelMenzel2009}). In the case of the present Slavic languages, four properties have turned out to characterise inflectional classes of nouns: \textit{Base form type}, \textit{stem type}, \textit{gender} and \textit{semantic animacy}.

As an example, let us consider first names. First names fall into two groups: male and female resp., which behave differently as to uninflectability. In the Slavic languages first names are characterised by the following semantic and morpho(no)logical properties: they are inherently animate and belong to either the masculine or the feminine gender, where a canonical male first name belongs to the masculine gender and a canonical female name to the feminine gender. Depending on the respective language, their base forms may either end in a vowel, which in Czech, Serbian, Slovak and Slovene may be phonologically short or long (while in Russian and Poish vowel length is not distinctive), or they may end in a consonant. Their stems may end in a vowel, a hard consonant, a soft consonant or a neutral (Czech, Slovak) sc. unpaired (Russian) consonant (cf. footnote 3).

It has now to be made clear, how a first name is assigned to an inflectional class, e.g., a borrowed name, such as the male name \textit{Pedro}. We can think of this assignment as a process of “feature checking”: First animacy and gender are established – in the case of \textit{Pedro}, masculine animate gender. Then the phonological ending of the name is matched with base forms of existing inflectional classes that contain masculine nouns. If one or several such classes exist, it has to be checked if the final vowel /o/ should be treated as part of the stem or as the Nominative ending for this particular class. If the latter is the case, the word form is analysed as \textit{Pedr-o} ‘STEM-ending’ and it has to be checked in a final step if the resulting stem conforms to the stem-structure of the chosen inflectional class. If it does, the word becomes a member of this class. If it does not, it remains uninflectable, unless it is adapted phonologically or morphologically, e.g., by changing the word ending and/or the stem ending. 

E.g., if the name \textit{Pedro} enters the Russian language, it is assigned to masculine-animate gender (which means that it selects a specific class of agreement targets). Afterwards it is checked, if there is an inflectional class with a base form ending in /o/ that allows for nouns with masculine-animate gender. It turns out that the only class with nouns ending in /o/ is reserved for neuter gender (\textit{okno} n. ‘window’), whereas the classes allowing for masculine-animate nouns require a base form ending in a consonant (\textit{Vladimir}) or in /a/ (\textit{Nikita}). Therefore, \textit{Pedro} cannot be assigned to any inflectional class and remains uninflectable in Russian.

If we take Slovene on the opposite end of the noun continuum in table 6, the same process yields a different result: In Slovene, there is one class specialized for masculine-animate nouns which at the same time allows for a base form in any vowel or consonant. In order to assign \textit{Pedro} to this class, the status of its stem has to be checked. The declension in question allows both for stems ending in a consonant and in a vowel. Since the final /o/ in \textit{Pedro} is unstressed, it is analysed as a NOM SG ending and the noun is inflected, like the native name \textit{Marko}, as \textit{Pedr-o} (NOM), \textit{Pedr-a} (GEN), \textit{Pedr-u} (DAT), \textit{Pedr-a} (ACC), (\textit{o}) \textit{Pedr-u} (LOC), (\textit{s}) \textit{Pedr-om} (INSTR). 

Tables 7 and 8 show the degree of uninflectability in first names in the six languages under investigation. This degree is again marked by different colourings of the cells – a darker colouring indicates a higher degree of uninflectability. A plus sign means that there are cases of systematic uninflectability for the base form ending in question, while a plus in brackets indicates cases of unsystematic uninflectability for this cell. White colouring combined with a minus sign means total lack of uninflectability. Shaded cells with question marks indicate that such cases could not be verified and probably do not exist.

As can be deduced from tables 7 and 8, uninflectability is much more widespread with female first names in all Slavic languages under investigation. As to male first names, only Russian and Polish display cases of uninflectability, while in Slovak, Czech, Slovene and Serbian all male first names can be assigned to an inflectional class, regardless of their phonological make-up.


\begin{tabularx}{\textwidth}{XXXXXXXXXXXXX}

\lsptoprule

male first names & {\bfseries base form ending in\footnotemark{}} & C\# & a & a: & o & o: & e & e: & i & i: & u & u:\\
& Russian & $-$ & \multicolumn{2}{c}{ +} & \multicolumn{2}{c}{ +} & \multicolumn{2}{c}{ +} & \multicolumn{2}{c}{ +} & \multicolumn{2}{c}{ +}\\
& Polish & $-$ & \multicolumn{2}{c}{ +} & \multicolumn{2}{c}{ (+)} & \multicolumn{2}{c}{ (+)} & \multicolumn{2}{c}{ (+)} & \multicolumn{2}{c}{ +}\\
& Slovak & $-$ & $-$ & $-$ & $-$ & $-$ & $-$ & $-$ & $-$ & $-$ & $-$ & $-$\\
& Czech & $-$ & $-$ & $-$ & $-$ & $-$ & $-$ & $-$ & $-$ & $-$ & $-$ & $-$\\
& Slovene & $-$ & $-$ & $-$ & $-$ & $-$ & $-$ & $-$ & $-$ & $-$ & $-$ & $-$\\
& Serbian & $-$ & $-$ & $-$ & $-$ & $-$ & $-$ & $-$ & $-$ & $-$ & $-$ & $-$\\
\lspbottomrule
\end{tabularx}
\footnotetext{ \textrm{In Polish and Russian vowel length is not distinctive, therefore the cells for the respective vowels are coalesced.}}
%%please move \begin{table} just above \\begin{tabular . 
\begin{table}
\caption{Uninflectability of male first names}
\label{tab:key:7}
\end{table}


\begin{tabularx}{\textwidth}{XXXXXXXXXXXXX}

\lsptoprule

female first names & {\bfseries base form ending in\textsuperscript{6}} & C\# & a & a: & o & o: & e & e: & i & i: & u & u:\\
& Russian & + & \multicolumn{2}{c}{ +} & \multicolumn{2}{c}{ +} & \multicolumn{2}{c}{ +} & \multicolumn{2}{c}{ +} & \multicolumn{2}{c}{ +}\\
& Polish & + & \multicolumn{2}{c}{ $-$} & \multicolumn{2}{c}{ +} & \multicolumn{2}{c}{ +} & \multicolumn{2}{c}{ +} & \multicolumn{2}{c}{ +}\\
& Slovak & + & $-$ & ? & + & + & + & + & + & + & + & +\\
& Czech & + & $-$ & ? & (+) & (+) & (+) & (+) & + & $-$ & + & +\\
& Slovene & (+) & $-$ & ? & (+) & (+) & (+) & (+) & + & + & + & +\\
& Serbian & (+) & $-$ & ? & (+) & (+) & (+) & (+) & + & + & + & +\\
\lspbottomrule
\end{tabularx}
%%please move \begin{table} just above \\begin{tabular . 
\begin{table}
\caption{Uninflectability of female first names}
\label{tab:key:8}
\end{table}

We will now concentrate on Russian first names and compare them to Slovene first names.

Russian inflectional classes are traditionally divided into three macroclasses\footnote{ \textrm{According to Dressler (1997, 2003) a class is a set of similar paradigms that can be subdivided into subclasses and ultimately into “microclasses which share exactly the same morphological generalizations. (…) An inflectional macroclass is the highest, most general type of class, which comprises several (sub)classes (…) Prototypically, its nucleus is a productive microclass and it has at least two microclasses. The interior coherence of a macroclass, in terms of shared properties, must be higher than affinities between microclasses of different macroclasses” \citep[35-36]{Dressler2003}.}} according to common sets of endings (cf. for Russian, e.g., \citealt{Timberlake2004}: 130). The 1\textsuperscript{st} macroclass displays (among others) a class containing nouns with masculine animate gender and a base form ending in a consonant, which is divided further into subclasses according to the stem-ending (identical to the base-form ending). Typical Russian male names belonging to this class are, e.g., \textit{Boris} (ending in a hard consonant), \textit{Igor’} (ending in a soft consonant), \textit{Jurij} (ending in an unpaired consonant). The 2\textsuperscript{nd} macroclass, on the other hand, requires a base form ending in /a/ and a stem ending in a consonant and can again be divided into subclasses depending on animacy. Animate nouns in this class may have feminine or masculine gender. Typical Russian female names here are, e.g., \textit{Anna} (stem ending in a hard consonant), \textit{Marija} (ending in a n unpaired consonant), (hypocoristic) \textit{Anja} (stem ending in a soft consonant), typical male names are \textit{Nikita} (stem ending in a hard consonant) and (hypocoristic) \textit{Kolja} (stem ending in a soft consonant). The 3\textsuperscript{rd} macroclass has a subclass with feminine animate gender and a base form (as well as stem) ending in a soft consonant which contains also traditional Russian female names, e.g., \textit{Ljubov’}, \textit{Raxil’} ‘Rachel’.

Let us first look at male names. As can be seen from table 7, in Russian male names ending in a vowel are affected by systematic uninflectability. But only in case of the vowels /e/, /i/ and /u/ is this correlation 100\%, i.e., it is the base form ending alone that conditions uninflectability. In other words, masculine-animate proper nouns with a base form ending in /e/, /i/ or /u/, do not conform to the conditions of any inflectional class in Russian. Therefore, foreign names as \textit{Arne, Teddi, Lu} are uninflectable: 

\ea%3
    \label{ex:key:3}
    \gll \\
         \\
    \glt
    \z % you might need an extra \z if this is the last of several subexamples

         \textit{Ja vstretila Arne/Teddi/Lu} 

I met:F Arne.ACC/Teddy.ACC/Lou.ACC

‘I met Arne/Teddy/Lou’.

With names ending in /o/ and /a/, other factors come into play, as can be seen in table 9. Names ending in /o/ are uninflectable, if the /o/ is stressed (\textit{Sadko}):

\ea%4
    \label{ex:key:4}
    \gll \\
         \\
    \glt
    \z % you might need an extra \z if this is the last of several subexamples

         \textit{Ja vstretila Sadko} 

I met:F Sadko.ACC’

‘I met Sadko’

Nevertheless, some traditional Slavic names ending in unstressed /o/ (\textit{Mi\textcyrillic{х}ajlo})\footnote{ \textrm{This form of the name “Michael” (nowadays} \textrm{\textit{Mixail}}\textrm{) was the original given name of the great Russian scholar Mixail Lomonosov. Such forms (as also} \textrm{\textit{Danilo}}, \textrm{\textit{Ivanko}}\textrm{) are obsolete as given names today. But they are still presented in dictionaries (Petrovskij n.y.) and characterised as inflectable according to the 2}\textrm{\textsuperscript{nd}} \textrm{declension which normally requires a base form ending in /a/, see Zaliznjak (2008: 739, also on} \textrm{\textit{Sadko}}).} may be inflected according to the 2\textsuperscript{nd} inflectional macroclass:

\ea%5
    \label{ex:key:5}
    \gll \\
         \\
    \glt
    \z % you might need an extra \z if this is the last of several subexamples

         \textit{Ja vstretila Mixajl-u} 

I met:F Mixajlo-ACC

‘I met Mixajlo’. 

This fact is reflected in table 9 by putting this base-form in brackets. (Unstressed /o/ is pronounced like unstressed /a/ as [ə], which may partly motivate this class assignment). 

In male names ending in /a/, uninflectability is conditioned by the stem ending. The corresponding nouns are assigned to the 2\textsuperscript{nd} declension, if /a/ is analysed as the Nominative suffix -a, as in \textit{Nikit-a}):

\ea%6
    \label{ex:key:6}
    \gll \\
         \\
    \glt
    \z % you might need an extra \z if this is the last of several subexamples

         \textit{Ja vstretila Nikit-u} 

I met:F Nikita-ACC

‘I met Nikita’. 

But membership in this class is allowed only if the remaining stem ends in a consonant. Therefore, a name such as (French) \textit{Fransua} “François” whose stem would end in /u/ (*Fransu-a) cannot be inflected:

\ea%7
    \label{ex:key:7}
    \gll \\
         \\
    \glt
    \z % you might need an extra \z if this is the last of several subexamples

         \textit{Ja vstretila Fransua} 

I met:F François.ACC 

‘I met François’. 

But (Japanese) \textit{Akira} is inflectable, since its stem ends in the consonant /r/ (Akir-a):

\ea%8
    \label{ex:key:8}
    \gll \\
         \\
    \glt
    \z % you might need an extra \z if this is the last of several subexamples

         \textit{Ja vstretila Akir-u} 

I met:F Akira-ACC

‘I met Akira’. 

Hence in Russian only male names ending in a consonant can be integrated into inflectional classes unconditionally, because in a concrete noun, all four constitutive properties – base form type, stem type, gender, subgender – have to have the fitting value. This implies that the constitutive properties of the Russian nominal inflectional classes are applied in a restrictive way and thus motivates the high degree of systematic uninflectability in male names.


\begin{tabularx}{\textwidth}{XXXXXX}

\lsptoprule

Class & base-form & gender & stem-type & Inflectable example & Uninflectable example\\
1. hard & C-Ø & m-anim. & hard C & \textit{Boris-Ø, Bob-Ø} & \textit{Pedro, Arne, Lu} “Lou”\\
~1. soft & C-Ø & m-anim. & soft C & \textit{Igor‘-Ø} & \textit{Mikele} „Michele“\textit{, Teddi} „Teddy“\\
1. unpaired\footnotemark{} & C-Ø & m-anim. & unpaired C & \textit{Jurij-Ø} & \textit{Mat‘je} „Mathieu“\textit{, Luidži} „Luigi“\\
2. & {}-a, 

(-o) & m-anim. & hard C & \textit{Nikit-a, (Mixajl-o)} & \textit{Fransua} „François“\textit{, Sadko (stressed /o/)}\\
~ & ~{}-a & m-anim. & soft C & \textit{Kolj-a} & \textit{Nikolja} »Nicolas« \textit{(stressed a and French origin)}\\
\lspbottomrule
\end{tabularx}
\footnotetext{ \textrm{In Russian /j/ belongs to the “unpaired consonants”.}}
%%please move \begin{table} just above \\begin{tabular . 
\begin{table}
\caption{Constitutive features of Russian masculine inflectional classes}
\label{tab:key:9}
\end{table}

The inflectional classes of Slovene are essentially analogous to those of Russian. Consequently, we will adhere to the numerical classification previously outlined.\footnote{ \textrm{Note that Slovene grammaticography defines macroclasses on the basis of gender, cf. \citet[40-61]{Herrity2000}, \citet[28-31]{Greenberg2008}.}} The 1\textsuperscript{st} macroclass displays again a class for nouns with masculine-animate gender. The 2\textsuperscript{nd} macroclass is characterised by a base form in /a/ and includes one class comprising feminine and masculine-animate nouns. (The 3\textsuperscript{rd} macroclass is reserved for feminine gender nouns, and is not productive in integrating first names).

With regard to male first names, these can belong to either the 1\textsuperscript{st} or the 2\textsuperscript{nd} macroclass (table 10). In the 1\textsuperscript{st} class, the base form of native nouns typically ends in a consonant (e.g., \textit{Borut}, \textit{Matjaž} ‚Matthew‘), in the 2\textsuperscript{nd}, it typically ends in /a/ (as in \textit{Miha}\footnote{ \textrm{Male names in unstressed /a/ can inflect either according to the 2}\textrm{\textsuperscript{nd}} \textrm{or the 1}\textrm{\textsuperscript{st}} \textrm{macroclass.}} ‘Michael’). Stems typically end in a consonant, as in Russian. However, in contrast to Russian, these inflectional classes permit a wide range of base forms and also allow for vocalic stems, as evidenced by the Latin name \textit{Leo}, where the final, unstressed /o/ is analysed as an ending and the noun inflected accordingly as \textit{Le-o} (NOM), \textit{Le-a} (GEN), \textit{Le-u} (DAT), \textit{Le-a} (ACC), (\textit{o}) \textit{Le-u} (LOC), (\textit{s}) \textit{Le-om} (INSTR). If the name ends in a stressed vowel or in /e/, /i/, /u/, the stem is automatically augmented by /j/\footnote{ \textrm{In Slovene /j/ belongs to the “soft consonants”.}} (or sometimes /t/) in the oblique cases, thereby creating a consonantal stem ending. This can be observed in the examples of \textit{Teddy-Ø, Teddyj-a} (GEN), \textit{Lou-Ø, Louj-a} (GEN), \textit{Tone-Ø}, \textit{Tonet-a} (GEN) ‘Tony’. Consequently, all male names are automatically assigned to an inflectional class. Therefore there are no uninflectable male names in Slovene.


\begin{tabularx}{\textwidth}{XXXXXX}

\lsptoprule

Class & base-form & gender & stem-type & inflectable & uninfl.\\
1. hard & C-Ø, 

{}-o, -a, -e (unstressed) & m-anim. & hard C or V & \textit{Borut-Ø, Mark-o, Mih-a}

\textit{Bob-Ø,} \textit{Pedr-o, Le-o}, \textit{Akir-a} & {}-\\
1. soft & C-Ø

V-Ø & m-anim. & soft C (stem-augmentation after vowel) & \textit{Matjaž-Ø}, \textit{Franz-Ø}

\textit{Nicolas-Ø stressed} /a/\textit{, François-Ø} stressed /a/ 

\textit{Arne-Ø, Teddy-Ø, Lou-Ø} & \textit{{}-}\\
2. & {}-a & m-anim. & C & \textit{Mih-a, Akir-a} & {}-\\
\lspbottomrule
\end{tabularx}
%%please move \begin{table} just above \\begin{tabular . 
\begin{table}
\caption{Constitutive features of Slovene masculine inflectional classes}
\label{tab:key:10}
\end{table}

We will now consider female first names. As previously noted, female names are significantly affected by uninflectability, as illustrated in \tabref{tab:key:8}. In Russian, all female names ending in a vowel other than /a/ are systematically uninflectable, e.g., \textit{Renate, Mimi, Ildiko, Lulu}. Furthermore, names with a base form in a consonant are also characterised by a high degree of uninflectability, e.g., \textit{Èdit} (German) “Edith”, \textit{Kler} “Claire”, and even certain names with a base form in /a/ may remain uninflectable, such as \textit{Klaudia} “Claudia”, \textit{Rea} “Rhea” (with two coda vowels).

In Slovene, uninflectability is also found in female first names, but systematically only with names ending in /i/ and /u/, such as \textit{Mimi}, \textit{Lulu}. Similarly, names with a base form ending in the vowels /e/ and /o/ as well as in a consonant are also affected by uninflectability, although not in a systematic manner: a corpus search for \textit{Renate Götschl}\footnote{ \textrm{A well-known Austrian athlete (1990’s-2000’s).}} in the Slovene corpus Gigafida yields 675 tokens for the word-form \textit{Renate}, and one token for the inflected form \textit{Renat-i} (DAT) (excluding items that may be related to a base form \textit{Renata}). The names \textit{Isabel} or \textit{Madeleine} usually appear in an uninflected form: A total of 3,123 tokens of the form \textit{Madeleine} are observed, along with one instance of the form \textit{Madelein-o} (ACC) and two instances of the form \textit{Madelein-i} (DAT). The same is true, mutatis mutandis, for \textit{Isabel}: there are 1,926 tokens of \textit{Isabel} and two of \textit{Isabel-i} (DAT) (where the relation to this base form, rather than an adapted \textit{Isabela} is undisputable). Such names are therefore classified as unsystematically uninflectable. In contrast, names ending in /a/ are unconditionally inflectable and go into the 2\textsuperscript{nd} (macro)class.

Let us again turn our attention to the inflectional classes of Russian and Slovene. In Russian, the conditions for female first names are similar to those for male first names: in order to be assigned to an inflectional class, a female name with feminine-animate gender has to provide a stem ending in a consonant and a base form ending in either /a/ or a soft consonant (the 3\textsuperscript{rd} macroclass). In reality, however, this latter case is obsolete, since the third declension is no longer productive (in the sense of \citealt{Dressler1997,Dressler2003}). This means that it is not open for borrowings or extragrammatical formations, even if they fulfil the conditions of the constitutive properties, e.g., the Chinese name \textit{Junin‘} “Yuning”. Nevertheless, older borrowings, such as biblical names (e.g., \textit{Raxil‘} “Rachel”), or the Soviet extragrammatical formation \textit{Ninel’} (an anagram of the name \textit{Lenin}) are inflectable (where \textit{Ninel‘} often remains uninflected, as a search of the Russian National Corpus\footnote{ \textrm{The research was conducted using the Russian National Corpus (}\href{https://ruscorpora.ru/}{{\textrm{ruscorpora.ru}}}\textrm{), \citealt{SavčukEtAl2024}.}} confirms). Therefore, this type is not classified as systematically uninflectable. New feminine\footnote{ \textrm{This has to be emphasized, because in Russian there are hypocoristic female names with masculine gender, which are derived by a native suffix and are inflected according to the 1}\textrm{\textsuperscript{st}} \textrm{macroclass, e.g.:} \textrm{\textit{Liza}} \textrm{f. >} \textrm{\textit{Liz-oček}} \textrm{m.}} female names can only be integrated into inflection if they have a base form that ends in /a/ and a stem that ends in a consonant. Therefore, names as \textit{Rea} (*Re-a), \textit{Klaudia} (*Klaudi-a) remain uninflectable. 

Female names can of course be adapted by changing an unfitting base form ending to /a/ (e.g., \textit{Renate} > \textit{Renata}) or by inserting the glide /j/ between two coda vowels (e.g., \textit{Klaudia} > \textit{Klaudija}).\footnote{\textrm{\citet{Ageenko2010} lists} \textrm{\textit{Klaudija Kardinale}} \textrm{(with glide-insertion) and} \textrm{\textit{Klaudia Schiffer}} \textrm{(without), \citet{Klyšinskij2010} recommends the form with /j/ for the transcription of Italian names in <ia> but without /j/ for German names with this ending.}} (This occurs frequently in casual speech and in the dialects). However, this is not a common practice in Contemporary Standard Russian, where usually unfitting names remain as they are and thus remain uninflectable.


\begin{tabularx}{\textwidth}{XXXXXX}

\lsptoprule

Macroclass & base-form & gender & stem-type & inflectable & uninflectable\\
2. hard & {}-a & f-anim. & hard C & \textit{Anna} & \textit{Rea} »Rhea«\\
2. soft & {}-a & f-anim. & soft C & \textit{Anja} & \\
3. & C-Ø & f-anim. & soft C, unpaired /š/, /ž/ & \textit{Raxil‘} & \textit{Solanž} »Solange«\textit{, Junin‘} „Yuning“\\
\lspbottomrule
\end{tabularx}
\textit{Table11: Constitutive features of Russian feminine inflection classes}

In Slovene the situation is more fluid, although the tendency is essentially the same: only names ending in /a/ are fully inflectable and are assigned to one of the subclasses of the 2\textsuperscript{nd} macroclass. In contrast to Russian, this class displays no restrictions as to the properties of animacy and stem-type. Therefore, both animate and inanimate nouns with a consonantal or vocalic stem can be assigned to it. Furthermore, names with a base form ending in /o/, /e/ or a consonant can be assigned to this class.  The part before such a vowel is analysed as the stem to which inflectional suffixes are attached. However, such cases are very rare and often historical, as with names from Greek mythology or from the Bible (\textit{Kli-o} (NOM), Kli\textit{{}-e} (GEN)). The 3\textsuperscript{rd} declension of Slovene, finally, does not contain or integrate female names at all.


\begin{tabularx}{\textwidth}{XXXXXX}

\lsptoprule

Class & base-form & gender & stem-type & inflectable & uninflectable\\
2. & {}-a, (o, e, C-Ø) & f & any & \textit{Alina, Lea,} (\textit{Klio, Melpomene, Ruth, Madeleine, Isabel}) & \textit{Lulu, Mimi, Ildiko}\\
3. & C-Ø & f & C & {}- & \\
\lspbottomrule
\end{tabularx}
\textit{Table12: Constitutive features of Slovene feminine inflection classes}

This discussion has shown that the nominal inflectional classes of Russian and Slovene are characterised by the constitutive properties \textit{gender} (including the subgender of \textit{animacy}), \textit{base form type} and \textit{stem type}. In these cases, the semantic category of \textit{animacy} is grammaticalised both as a subgender and a subclass of the corresponding inflectional classes. 

However, animacy may also be a semantic constitutive property, as evidenced by Czech and Slovak. In Czech, there is one inflectional class for feminine nouns ending in a soft consonant, analogous to the respective 3rd class in Russian and Slovene. This class is productive, but it does not contain animate nouns. Because of this semantic constitutive property, it integrates common nouns such as \textit{kartridž-Ø} ‘cartridge’, toponyms as \textit{Marseille} (ending in /j/) and \textit{Provence} (ending in /s/), but a female name such as \textit{Florence} (ending in /s/) cannot be assigned to this class and remains uninflectable. 

\section{Summary} %5 /

The objective of this article is to demonstrate the following points: Uninflectedness in inflecting languages is not a uniform phenomenon. In the first place, two fundamentally distinct manifestations can be identified: contextual and paradigmatic uninflectedness. The former represents a temporary absence of inflection in an otherwise inflectable lexeme, whereas the latter is a permanent feature of either a single lexeme or a class of lexemes with identical properties. The first type has been termed \textit{contextual uninflectedness}, while the two varieties of the second type have been termed \textit{systematic} and \textit{unsystematic uninflectability}. These manifestations can be attributed to different causes which have been spelled out for the noun.

Contextual uninflectedness is caused by contextual factors, such as syntactic context or register, but may also be sociolinguistically conditioned. 

Systematic uninflectability of nouns is caused either by a mismatch between the constitutive properties of inflectional classes and the properties of the nouns in question, or by the fact that an otherwise matching class is unproductive.

Finally, unsystematic uninflectability can on the one hand occur for reasons of style or register. On the other hand, it can be an indicator of ongoing changes, showing either that an inflectional class is changing its constitutive properties or that it is becoming unproductive. 

\subsection{References}

Ageenko, Florencija L. 2010. Slovar’ sobstvennyx imen russkogo jazyka. Udarenie. Proiznošenie. Slovoizmenenie. Bolee 38 000 slovarnyx edinic. Moskva: Mir i obrazovanie. \url{https://gramota.ru/biblioteka/slovari/slovar-sobstvennyh-imyon} 

Brown, Dunstan;  Marina  Chumakina;  Greville  Corbett;  Gergana  Popova  \& Andrew Spencer. 2012. Defining ‘periphrasis’: key notions. Morphology.  DOI 10.1007/s11525-Ø12-9201-5 

Comrie, Bernard, Corbett, Greville. 1993. Introduction. In: Comrie, Bernard, Corbett, Greville (eds.). The Slavonic Languages. London: Routledge, 1-19.

Doleschal, Ursula. 1997. Indeklinabilität, Semideklinabilität und Flexionsklassen: Russisch und Slowakisch. Wiener Slavistisches Jahrbuch 43, 41-52.

\_\_\_\_\_. 1999. „Milka“, „Duracell“ und andere adjektivische Wortungereimtheiten In: Anstatt, Tanja, Roland Meyer \& Elisabeth Seitz (Hg.), Linguistische Beiträge zur Slavistik aus Deutschland und Österreich. München: Sagner, 87-98

\_\_\_\_. 2000. \textit{Das Phänomen der Unflektierbarkeit in den slawischen Sprachen}. Unpubl. habil. thesis, Wirtschaftsuniv. Wien.

\_\_\_\_. 2002. „Super sexy Mici“ oder das Phänomen der Unflektierbarkeit in den slavischen Sprachen. In: Daiber, Thomas (Hg.), \textit{Linguistische Beiträge zur Slavistik aus Deutschland und Österreich}. München: Sagner, 55-66.

\_\_\_\_. 2006. Konstitutive Eigenschaften von Flexionsklassen. In: E. Binder, W. Stadler, H. Weinberger (Hg.). \textit{Zeit - Ort - Erinnerung. Slawistische Erkundungen aus sprach-, literatur- und kulturwissenschaftlicher Perspektive}. Innsbruck: Inst. F. Sprachwissenschaft, 379-394

\begin{styleBibliographie}
\_\_\_\_\_. 2007. Über die Entstehung des Deklinationstyps kuli, kuliho. In: Geist Ljudmila, Mehlhorn, Grit (Hg.). Linguistische Beiträge zur Slavistik, XIV. JungslavistInnen-Treffen in Stuttgart. München: Sagner, 51-64
\end{styleBibliographie}

\_\_\_\_\_. 2009. Nominale Kategorien: Genus. In: Berger, Tilman, Gutschmidt, Karl, Kempgen, Sebastian, Kosta, Peter (eds). Die slavischen Sprachen / The Slavic Languages: Ein internationales Handbuch zu ihrer Struktur, ihrer Geschichte und ihrer Erforschung / An International Handbook of their Structure, their History and their Investigation. Halbband 1. Berlin: Mouton – De Gruyter. 143-152

\_\_\_\_\_. 2016. „Završena velika Konzum nagradna igra“ – on the status of premodifying nouns in Croatian. \textit{Fluminensia} 27, 2015/2, 191–202. hrcak.srce.hr/file/223577

\_\_\_\_\_. 2017. „SMS nagradna igra. Eine Konstruktion zwischen Syntax und Wortbildung“. In: \textit{Wiener Slawistischer Almanach, Band 79 (2017), 001-Ø07}

Dressler, Wolfgang U. 1997. On Productivity and Potentiality in Inflectional Morphology. CLASnet Working Papers 7 (Montréal), 3-22.

\_\_\_\_. 2000. Extragrammatical vs. Marginal Morphology. In: Doleschal, Ursula \& Anna Maria Thornton (eds.): Extragrammatical and Marginal Morphology. München: LINCOM Europa, 1-10.

\_\_\_\_. 2003. Degrees of grammatical productivity in inflectional morphology. In: Rivista di Linguistica 15.1, 31-62.

Dressler, Wolfgang U., Katarzyna Dziubalska-Kołaczyk \& Małgorzata Fabiszak. 1997. Polish inflection classes within Natural Morphology. Bulletin de la Société polonaise de linguistique 53, 95-119. 

Genis, René. 2021. “Aspectual Pairs”, in: Encyclopedia of Slavic Languages and Linguistics Online, Editor-in-Chief General Editor Marc L. Greenberg, Lenore A. Grenoble. Consulted online on 21 \citealt{November2023} <http://dx.doi.org/10.1163/2589-6229\_ESLO\_COM\_038613> First published online: 2021

Gigafida 2.0: Korpus pisne standardne slovenščine, viri.cjvt.si/gigafida, dostop 24. 6. 2024.

Graudina, Ljudmila K. 1980. Voprosy normalizacii russkogo jazyka. Grammatika i varianty. Moskva: Nauka.

Greenberg, \citealt{Marc2008}. A short reference grammar of Slovene. München: LINCOM Europa.

\begin{styleLiteraturangabe}
Haspelmath, Martin. 2000. Periphrasis. In: Booij, Geert, Christian Lehmann~\& ~Joachim Mugdan (eds). Morphology: A Handbook on Inflection and Word Formation. Vol.~1. Berlin: de Gruyter, 654–664.
\end{styleLiteraturangabe}

Hentschel, Gerd, Thomas Menzel. 2009. „Nominale Kategorien: Kasus.“ In: Berger, Tilman, Gutschmidt, Karl, Kempgen, Sebastian, Kosta, Peter (eds). Die slavischen Sprachen / The Slavic Languages: Ein internationales Handbuch zu ihrer Struktur, ihrer Geschichte und ihrer Erforschung / An International Handbook of their Structure, their History and their Investigation. Halbband 1. Berlin: Mouton – De Gruyter. 161-176.

Herrity, Peter. 2000. Slovene: A Comprehensive Grammar. London: Routledge

Internetová jazyková příručka [online] (2008–2024). \href{https://prirucka.ujc.cas.cz/#div398}{{Skloňování názvů institucí a firem}}{.} Praha: Ústav pro jazyk český AV ČR, v.\hspace{0.2em}v.\hspace{0.2em}i. Cit. 11.\hspace{0.2em}5.\hspace{0.2em}2024. <\url{https://prirucka.ujc.cas.cz/}>. \url{https://prirucka.ujc.cas.cz/?id=398} 

Internetová jazyková příručka [online] (2008–2024). Zkratková slova. Praha: Ústav pro jazyk český AV ČR, v.\hspace{0.2em}v.\hspace{0.2em}i. Cit. 11.\hspace{0.2em}5.\hspace{0.2em}2024. <\url{https://prirucka.ujc.cas.cz/}>. \url{https://prirucka.ujc.cas.cz/?id=784} 

Klyšinskij, È.~S. (ed.). 2010. Praktičeskaja transkripcija ličnych imen v jazykach narodov mira. Moskva: Nauka.

Nagórko, Alicja. 2009. Struktur des Lexikons. In: Berger, Tilman, Gutschmidt, Karl, Kempgen, Sebastian, Kosta, Peter (eds). Die slavischen Sprachen / The Slavic Languages: Ein internationales Handbuch zu ihrer Struktur, ihrer Geschichte und ihrer Erforschung / An International Handbook of their Structure, their History and their Investigation. Halbband 1. Berlin: Mouton – De Gruyter.739-758.

Orzechowska, Alicja. 1974. Niedomienne rzeczowniki pospolite w języku polskim. Studia z filologii polskiej i słowiańskiej 14, 103-114.

Paxomov, V.M. n.y.  \textcyrillic{В Простоквашино или в Простоквашине?} Azbučnye istiny. Gramota.ru. \url{https://gramota.ru/biblioteka/spravochniki/azbuchnye-istiny/istiny-1-toponimy} (19.05.2024)

Panov, Michail V. (red.). 1968. Russkij jazyk i sovetskoe obščestvo. Morfologija i sintaksis sovremennogo russkogo literaturnogo jazyka. Moskva: Nauka.

Petrovksij, N.A. no year. Slovar’ russkix imen. \url{https://gramota.ru/biblioteka/slovari/slovar-russkih-imyon} (23.06.2024)

Považaj, Matej. Red. 2013. Pravidlá slovenského pravopisu. 4., nezmenené vydanie. Bratislava: Veda. \url{https://www.juls.savba.sk/psp_2013.html} 

\_\_\_\_\_\_\_. 1998. Pravidlá slovenského pravopisu. 2. dopln. a preprac. vydanie. Bratislava: \citealt{Veda1998}. 

Redakcija portala «Gramota.ru» n.y. Kak sklonjat‘ geogafičeskie nazvanija? \url{https://gramota.ru/biblioteka/spravochniki/pismovnik/kak-sklonyat-geograficheskie-nazvaniya} (19.05.2024)

Russian National Corpus: \textcyrillic{Савчук С. О., Архангельский Т. А., Бонч-Осмоловская А. А., Донина О. В., Кузнецова Ю. Н., Ляшевская О. Н., Орехов Б. В., Подрядчикова М. В. Национальный корпус русского языка 2.0: новые возможности и перспективы развития. Вопросы языкознания, 2024, 2: 7–34.}

Słownik języka polskiego. (red.,), \emph{Słownik języka polskiego}, t. 1–11, \citealt{Warszawa1958}–1969, PWN. \url{https://sjp.pwn.pl/} 

Spencer, A. 2020. Uninflectedness: Uninflecting, Uninflectable and Uninflected Words, or the Complexity of the Simplex. In Livia Körtvélyessy \& Pavol Štekauer (Eds.), Complex Words: Advances in Morphology (pp. 142-158).

Spencer, Andrew and Wiemer, Björn. 2020. “Inflectional Systems”, in: Encyclopedia of Slavic Languages and Linguistics Online, Editor-in-Chief General Editor Marc L. Greenberg, Lenore A. Grenoble. Consulted online on 21 \citealt{November2023} <http://dx.doi.org/10.1163/2589-6229\_ESLO\_COM\_033721>.First published online: 2020

Suprun, Adam E. 1969. Slavjanskie čisliteľnye. Stanovlenie čisliteľnych kak osoboj časti reči. Minsk.

Sussex, Roland; Cubberley, Paul. 2006: The Slavic Languages. Cambridge: Cambridge Univ. Press.

Timberlake, Alan. 2004. A Reference Grammar of Russian.Cambridge: CUP.

Tejkin, M.S. 2021. K voprosu o sklonenii toponimov srednego roda na \textit{{}-ovo/-evo, -ino/-yno}. Vestnik Kemerovskogo gosudarstvennogo universiteta 23/4: 1074-1085. \url{https://doi.org/10.21603/2078-8975-2021-23-4-1074-1085} 

Woliński, Marcin et al. \textsuperscript{4}2020. Słownik gramatyczny języka polskiego. Warszawa \url{http://sgjp.pl/o-slowniku/}Wurzel, Wolfgang U. 1984. Flexionsmorphologie und Natürlichkeit. Berlin: Akademie-Verlag.

\begin{styleLiteraturangabe}
Worth, Dean S. 1966. On the Stem/ending Boundary in Slavic Indeclinables. Zbornik za filologiju i lingvistiku 9, 10–16.
\end{styleLiteraturangabe}

Zaliznjak, Andrej Anatolʹevič, 2008. Grammatičeskij slovar' russkogo jazyka: slovoizmenenie. Izd. 5-e, ispr. Moskva: AST-press.

Zgusta, Ladislav. 1996. Names and their study. In: Eichler, Ernst, Gerold Hilty, Heinrich Löffler, Hugo Steger und Ladislav Zgusta (eds.). Namenforschung / Name Studies / Les noms propres. Ein internationales Handbuch zur Onomastik / An International Handbook of Onomastics / Manuel international d'onomastique. 2. Halbband+Registerband. Berlin, Boston: De Gruyter Mouton. 1876-1890.\url{https://doi.org/10.1515/9783110203431}. 

Żmigrodzki, Piotr (red.) 2004-2021. Wielki słownik języka polskiego. Instytut Języka Polskiego PAN. \url{https://wsjp.pl/}


\begin{verbatim}%%move bib entries to  localbibliography.bib
\end{verbatim}
\sloppy\printbibliography[heading=subbibliography,notkeyword=this]
\end{document}
